% ----------------------------------------------------------
% Style for submissions to Nature Portfolio journals
% ----------------------------------------------------------
\documentclass[pdflatex,sn-nature]{sn-jnl}
% ----------------------------------------------------------
% additional latex packages 
% ----------------------------------------------------------
\usepackage{graphicx}
\usepackage{multirow}
\usepackage{amsmath,amssymb,amsfonts}
\usepackage{amsthm}
\usepackage{mathrsfs}
\usepackage{xcolor}
\usepackage[title]{appendix}%
\usepackage{textcomp}
\usepackage{manyfoot}
\usepackage{booktabs}
\usepackage{algorithm}
\usepackage{algorithmicx}
\usepackage{algpseudocode}
\usepackage{listings}

\raggedbottom
% ----------------------------------------------------------https://www.overleaf.com/project/641981d1f29991b6c00119e5
% Inicio do Artigo ChatGPT e Design Patterns
% ----------------------------------------------------------
\begin{document}
\title{Can ChatGPT suggest design patterns? A study based on the analysis of exam questions.}
% ----------------------------------------------------------
% Nome autores
% ----------------------------------------------------------
\author[1]{\pfx{Dr.} \fnm{Eduardo} \spfx{Martins} \sur{Guerra}}\email{guerraem@gmail.com}

\author[2]{\fnm{Filipe}  \sur{Correia}} \email{filipe.correia@fe.up.pt}

\author[3]{\fnm{João} \spfx{José} \sur{Maranhão}  \sfx{Junior}}\email{jjjunior@gmail.com}

% ----------------------------------------------------------
% Filiação
% ----------------------------------------------------------
\affil[1]{\orgdiv{Computer Science}, \orgname{Free University of Bozen-Bolzano}, \orgaddress{\street{Piazza Università, 1}, \city{Bolzano}, \postcode{39100}, \state{BZ}, \country{Italy}}}

\affil[2]{\orgdiv{Department}, \orgname{Organization}, \orgaddress{\street{Street}, \city{City}, \postcode{10587}, \state{State}, \country{Country}}}

\affil[3]{\orgdiv{Computação Aplicada}, \orgname{IPT - Instituto de Pesquisas Tecnologicias}, \orgaddress{\street{Av. Prof. Almeida Prado, 532}, \city{São Paulo}, \postcode{05508-901}, \state{SP}, \country{Brazil}}}

% ----------------------------------------------------------
% Abstract
% ----------------------------------------------------------
\abstract{This study aims to evaluate the effectiveness of ChatGPT in suggesting appropriate design patterns to solve software development problems. To achieve this, we collected 53 pre-existing questions from various Internet sources with scenario descriptions and asked ChatGPT to interpret and respond. We then assessed the quality of the ChatGPT responses against pre-existing question feedback to determine each suggestion's accuracy, relevance, and comprehensiveness. The results of this study provide insight into the potential of AI-assisted tools for software development and the role they can play in improving software quality.}

\keywords{ChatGPT, Design Patterns, Artificial Intelligence, I-assisted}

\maketitle

\section{Introduction}\label{sec1}
The use of design patterns\cite{bib1,bib3} in software development is prevalent in software engineering. It brings several advantages, including code reuse, better code organization, ease of maintenance, better scalability, and software quality\cite{bib4}. However, becoming proficient in understanding and applying these patterns can take time and effort\cite{bib2}, requiring a high degree of abstraction and practice in object-oriented programming. Educational institutions use various learning techniques\cite{bib5} to engage students in real-world problems to recognize and interpret these patterns.

ChatGPT is an extensive language model trained on a large text dataset that includes much information about software development, including design patterns. As a result, these AI-assisted tools can suggest design patterns based on a specific set of requirements or a specific problem. To indicate a design pattern, he would need to be given information about the specific problem that needs to be addressed and any relevant constraints or requirements. This can be done through a conversation or by providing a written problem description. Once this artificial intelligence clearly understands the problem, it can use its knowledge of software design patterns to suggest the most appropriate design pattern. This would involve analyzing the situation, identifying essential requirements, and considering the pros and cons of each design pattern.

This study aims to evaluate the effectiveness of ChatGPT responses to assist developers in choosing a design pattern to solve a design scenario for this purpose; We collected 53 pre-existing questions from various Internet sources with scenario descriptions and asked them to interpret and respond.

% Que aprender quando aplicar padrões não é simples//ok
% ChatGPT//ok
% Possiblidade de AI ajudar
	.	
% o que a gente fez//ok
	
% resumo dos resultados
	
% descreve seções



\section{Background}\label{sec2}

\subsection{AI-Assisted Tools for Implementing Design Patterns}\label{sec2subsec1}

\subsection{AI-Based knowledge extraction for automatic design proposals using design-related patterns}\label{sec2subsec2}

\subsection{Towards Predicting Architectural Design Patterns: A Machine Learning Approach
}\label{sec2subsec3} 

\section{Teaching Design Patterns}\label{sec3}

\subsection{A problem-based approach to teaching design patterns}\label{sec3subsec1}
\subsection{Study Design Patterns}\label{sec3subsec2}

\section{Research Questions}\label{sec4}
\textbf{How effective is ChatGPT in suggesting appropriate design patterns to solve software development problems?}\newline

 This research question will guide the study to investigate the accuracy, relevance, and comprehensiveness of ChatGPT responses by suggesting design patterns for the 53 pre-existing questions asked. The results of this study will contribute to our understanding of the potential of AI-assisted tools in software development and their role in improving software quality.

\section{Research Methodology}\label{sec5}
In the initial phase of our research, we endeavored to identify the type of design pattern best suited for our study. We chose the catalog of 23 design patterns from the famous book “Gang of Four” (GOF); we selected the questions that make up the standards from this catalog that are widely disseminated design patterns in the job market and extensively taught in higher education or open courses. 

In the next phase, we looked for pre-existing questions prepared by higher education professors or public tenders.
We then filtered out the most pertinent questions that provided a clear context or explanation of an issue without directly referencing the pattern. 

In the final phase, we entered 53 questions into ChatGPT to determine the most appropriate design pattern for each question. We then validated whether the recommended pattern was correct or incorrect. Our dataset comprised 20 questions in English and 33 questions in Portuguese.

\subsection{What is the GOF design pattern?}\label{sec5subsec1}
The GoF design pattern refers to the patterns presented in the book "Design Patterns: Elements of Reusable Object-Oriented Software" written by Erich Gamma, Richard Helm, Ralph Johnson, and John Vlissides, also known as the "Gang of Four" (GOF). The book describes a set of common design patterns for solving recurring problems in object-oriented software development.
% -Categorizar os padroes
% -Explicar cada padrão

\subsection{Questions Dataset}\label{sec5subsec2}
    % visão geral das questões: por categoria de padrão, quantidade, por padrão, por fonte, lingua, etc...


\subsection{Questions Analysis}\label{sec5subsec3}
We used a mixed methods approach to analyze the questions collected in this study. First, we evaluate ChatGPT responses against known appropriate design patterns for each question. We assess each suggestion's accuracy, relevance, and completeness to determine the ChatGPT response's quality. We then analyzed the data statistically to identify significant differences between the ChatGPT suggestions and known design patterns.

\subsection{Validity and Reliability}\label{sec5subsec4}
To ensure the validity and reliability of the study, we used a rigorous data analysis process and compared ChatGPT suggestions against known and appropriate design standards for each question.

% passos do estudo
% GoF
%selecionar questões (fala do nosso critério)
% avaliou a resposta: métricas - taxa de acerto, avaliação qualitativa do texto, de quais padrões ele errou e acertou

\subsection{Participants}\label{sec5subsec5}
No human participants were involved in this study, as we used pre-existing questions from internet sources.

\subsection{Ethics}\label{sec5subsec6}
This study did not involve human subjects, so no ethical approval was required. We ensured that the pre-existing questions used in the study were publicly available and did not contain confidential or sensitive information.


\section{Results}\label{sec6}

	% 4.1 seco sem interpretação
	
	% 4.2 quantidade de acerto
	
	% 4.3 para os erros, os detalhes da questão: por padrão, por idioma etc...
	
	% 4.4 análise do tamanho das respostas (média, min e max nr de chars)

\section{Discussion}\label{sec7}
 \begin{enumerate}
     \item \textbf{How is the success rate of ChatGPT in the design patterns exercises?}
     \item \textbf{How does the answer of ChatGPT help in the pattern implementation?}
     \item \textbf{In questions where the wrong pattern is pointed, does ChatGPT answer mislead the developers?}
 \end{enumerate}
 
\section{Conclusion}\label{sec8}
    % 	Resumo
    % 	Principais resultados
    % 	Trabalhos futuros
    % 	Outros padrões
    % 	Problemas da Industria

\bibliography{references}

\end{document}
